\paragraph{Fermion replica gluing for \(\operatorname{Tr}\rho^n\) and the supertrace character (conditional; not closure).}
Let the state of the Poisson-limit fermions on the wedge be quasi-free, and let its modular flow be the rotation by \(2\pi t\) of the transverse plane of \(\Sigma\), acting on every field through its representation of the rotation group (for integer spin, the identity at integer \(n\)). The ungraded trace \(\operatorname{Tr}\rho^n\) then glues a Grassmann-odd field with the rotation by \(2\pi n\) and one sign from the cyclicity of the trace: in the polar frame the spinor is antiperiodic under the angle \(2\pi n\), and the Cartesian components of transverse spin \(\sigma=\pm\tfrac12\) acquire the monodromy \(-e^{2\pi\mathrm in\sigma}\), which is \((-1)^{n+1}\) for integer \(n\)~\cite{CasiniHuerta2009,FursaevMiele1997}. Diagonalising the replica matrix gives the twist phases \(e^{2\pi\mathrm ik/n}\), \(k=-(n-1)/2,\dots,(n-1)/2\)~\cite{CasiniHuerta2009}. For a block of two components with twist \(a\), \(|a|<\tfrac12\), and the less singular self-adjoint extension of the first-order operator \(\gamma^\mu D_\mu\), whose square is \(\mathcal O_f\), the conical heat-trace constant is \(a^2+\tfrac16\): the componentwise-regular value \(B_2(|a|)\), \(B_2(x)=x^2-x+\tfrac16\), plus \(|a|\) from the reversed sign of the one Bessel order that the first-order operator fixes. This extension is the one giving an analytic continuation: the componentwise-regular domain gives \(\sum_k|a_k|\) equal to \(n/4\) or \((n^2-1)/(4n)\) for even or odd \(n\), a parity-dependent and non-analytic \(\operatorname{Tr}\rho^n\). With the less singular extension, \(\sum_ka_k^2=(n^2-1)/(12n)\), analytic in \(n\), equal to the zeta-function value of the polar-frame cone on the branch \(n<1\), with \((1-n\partial_n)\) equal to \(-\tfrac16\) per block at \(n=1\). The area coefficient of a Dirac field (two blocks, weight \(-1\)) is therefore \(+\tfrac13\), as used above for antiperiodic gluing.

The unrotated periodic gluing above has the phase \(+1\) for all \(n\): its twists are those of \(\operatorname{Tr}\rho^n\) for odd \(n\) and those of the graded trace \(\operatorname{Tr}[(-1)^F\rho^n]\) for even \(n\), and with the componentwise-regular domain it gives \(-2\) times the Dirac-domain value at odd \(n\). It is not \(\operatorname{Tr}\rho^n\), and for the stated extension its values \(-\tfrac23\) and \(4-2n_D\) do not apply to \(\operatorname{Tr}\rho^n\). The graded trace has the twist \(a=\tfrac12\) at \(n=1\) and is not normalised there. The rotation by \(2\pi n\) is the identity on the vector index for integer \(n\), so the integer-\(n\) traces of the connection sector are those of the unrotated gluing; if the continuation from integer \(n\) is unique, the rotated branches \(3-n_D\) and \(5-n_D\) above do not reproduce them and are not continuations of \(\operatorname{Tr}\rho^n\). With \(c_A=\tfrac{N_s+\alpha+1}6+\tfrac{n_D}3\) and \(c_T\) as above, \(c_A-c_T=4(1-\tfrac{n_D}2)+n_D\), the sum of the weight of \(\rho^2X\) and the excess \(+1\) per Dirac field of the cone value \(+\tfrac13\) over its coefficient \(-\tfrac23\) of \(\rho R[g_\rho]\). For \(n_D=2\), \(N_s=6\), \(\alpha=1\): \(c_T=0\), \(c_A=2\).

In the character \(\mathsf Q_{10}\) of~\cite{MantaSteinackerTran2024}, Eq.~(3.7), with the flux along the four Poisson directions, the vector, six inert scalars, the ghosts and \(n_D\) Dirac fields in place of the sixteen-component spinor, the coefficient of order \(\alpha^0\) is \(N_{\mathrm{tot}}=8-4n_D\) and that of order \(\alpha^2\) is \(-b\,(x^2+y^2)\) with \(b=1-n_D/2\), where \(\pm\mathrm ix,\pm\mathrm iy\) are the eigenvalues of the flux. These are the weighted component count and the weighted index of the preceding paragraphs, so that \(c_T=N_{\mathrm{tot}}/6-4b\) vanishes at \(n_D=2\), \(N_s=6\) because the orders \(\alpha^0\) and \(\alpha^2\) of the character vanish. At \(n_D=2\) the order \(\alpha^4\) is \((x^2-y^2)^2/16\), which vanishes for self-dual or anti-self-dual flux. This concerns the character only, not the \(a_2\) coefficient of the operators and not the area side. These statements are conditional on the quasi-free state, the rotational modular flow, the choice of self-adjoint extension, the uniqueness of the continuation from integer \(n\), the domain of the first-order operator, the \(\LUV^2\) order and the class of the preceding paragraphs; they do not derive the state, the flow, the content, or a value or sign of \(G_{\mathrm{eff}}\), and they do not compute the next order of either side. Problems~B, C, D, F and~G remain open.
